\begin{ExerciseBlock}
\ExerciseHeading{}{Verständnisfragen}{forced_oscillations_checks}
Entscheiden Sie, welche Aussagen zutreffen, und begründen Sie Ihre Entscheidung.
\begin{enumerate}
\item Die Frequenz einer erzwungenen Schwingung ist im eingeschwungenen Zustand kleiner als die Erregerfrequenz.\TrueFalseAnswer{falsch}
\item Bei sonst gleichen Parametern ist das Amplitudenmaximum für kleine Dämpfungen besonders hoch.\TrueFalseAnswer{wahr}
\item Weit unterhalb der ungedämpften Eigenfrequenz sind Auslenkung und Erregerkraft näherungsweise in Phase.\TrueFalseAnswer{wahr}
\item Bei sehr starker Dämpfung gibt es kein Amplitudenmaximum bei einer Erregerfrequenz größer als null.\TrueFalseAnswer{wahr}
\end{enumerate}
\end{ExerciseBlock}
